\SourceLeaf{092}

{\rightskip=0pt plus 1em Aggregatum quadratorum nostrarum compensationum invenitur \(=97{,}8845\). Hinc error medius, quatenus ex 27 angulis observatis colligi potest,\par}
\[
=\BareRoot{\frac{97{,}8845}{13}}=2''7440
\]

Aggregatum quadratorum mutationum, quas clar. \textsc{de Krayenhof} ipse angulis observatis applicavit, invenitur \(=341{,}4201\).

\Article{24}
Exemplum alterum suppeditabunt triangula inter quinque puncta triangulationis Hannoveranae, Falkenberg, Breithorn, Hauselberg, Wulfsode, Wilsede. Observatae sunt directiones\footnote{Initia, ad quae singulae directiones referuntur, hic tamquam arbitraria considerantur, quamquam revera cum lineis meridianis stationum coincidunt. Observationes in posterum complete publici iuris fient; interim figura invenitur in \textit{Astronomische Nachrichten} Vol. I. p. 441.}:
\begin{center}
\begin{tabular}{r l l r r r}
\multicolumn{6}{c}{In statione \textit{Falkenberg}}\\[5pt]
0.&Wilsede&. . . . .&\(187^0\)&\(47'\)&\(30''311\)\\
1.&Wulfsode&. . . .&225&9&39,676\\
2.&Hauselberg&. . .&266&13&56,239\\
3.&Breithorn&. . . .&274&14&43,634
\end{tabular}
\end{center}
\begin{center}
\begin{tabular}{r l l r r r}
\multicolumn{6}{c}{In statione \textit{Breithorn}}\\[5pt]
4.&Falkenberg&. . .&94&33&40,755\\
5.&Hauselberg&. . .&122&51&23,054\\
6.&Wilsede&. . . . .&150&18&35,100
\end{tabular}
\end{center}
\begin{center}
\begin{tabular}{r l l r r r}
\multicolumn{6}{c}{In statione \textit{Hauselberg}}\\[5pt]
7.&Falkenberg&. . .&86&29&6,872\\
8.&Wilsede&. . . . .&154&37&9,624\\
9.&Wulfsode&. . . .&189&2&56,376\\
10.&Breithorn&. . . .&302&47&37,732
\end{tabular}
\end{center}
\begin{center}
\begin{tabular}{r l l r r r}
\multicolumn{6}{c}{In statione \textit{Wulfsode}}\\[5pt]
11.&Hauselberg&. . .&9&5&36,593\\
12.&Falkenberg&. . .&45&27&33,556\\
13.&Wilsede&. . . . .&118&44&13,159
\end{tabular}
\end{center}
\SourceLeaf{093}
\begin{center}
\begin{tabular}{r l l r r r}
\multicolumn{6}{c}{In statione \textit{Wilsede}}\\[5pt]
14.&Falkenberg&. . .&\(7^0\)&\(51'\)&\(1''027\)\\
15.&Wulfsode&. . . .&298&29&49,519\\
16.&Breithorn&. . . .&330&3&7,392\\
17.&Hauselberg&. . .&334&25&26,746
\end{tabular}
\end{center}

Ex his observationibus septem triangula formare licet.
\begin{center}
\begin{tabular}{l l r r r}
\multicolumn{5}{c}{Triangulum I.}\\[5pt]
Falkenberg&. . . . . .&\(8^0\)&\(0'\)&\(47''395\)\\
Breithorn&. . . . . . .&28&17&42,299\\
Hauselberg&. . . . . .&143&41&29,140
\end{tabular}
\end{center}
\begin{center}
\begin{tabular}{l l r r r}
\multicolumn{5}{c}{Triangulum II.}\\[5pt]
Falkenberg&. . . . . .&86&27&13,323\\
Breithorn&. . . . . . .&55&44&54,345\\
Wilsede&. . . . . . . .&37&47&53.635
\end{tabular}
\end{center}
\begin{center}
\begin{tabular}{l l r r r}
\multicolumn{5}{c}{Triangulum III.}\\[5pt]
Falkenberg&. . . . . .&41&4&16,563\\
Hauselberg&. . . . . .&102&33&49,504\\
Wulfsode&. . . . . . .&36&21&56,963
\end{tabular}
\end{center}
\begin{center}
\begin{tabular}{l l r r r}
\multicolumn{5}{c}{Triangulum IV.}\\[5pt]
Falkenberg&. . . . . .&78&26&25,928\\
Hauselberg&. . . . . .&68&8&2,752\\
Wilsede&. . . . . . . .&35&25&34,281
\end{tabular}
\end{center}
\begin{center}
\begin{tabular}{l l r r r}
\multicolumn{5}{c}{Triangulum V.}\\[5pt]
Falkenberg&. . . . . .&37&22&9,365\\
Wulfsode&. . . . . . .&73&16&39,603\\
Wilsede&. . . . . . . .&69&21&11,508
\end{tabular}
\end{center}
\begin{center}
\begin{tabular}{l l r r r}
\multicolumn{5}{c}{Triangulum VI.}\\[5pt]
Breithorn&. . . . . . .&27&27&12,046\\
Hauselberg&. . . . \hspace{.5em} .&148&10&28,108\\
Wilsede&. . . . . . . .&4&22&19,354
\end{tabular}
\end{center}
\SourceLeaf{094}
\begin{center}
\begin{tabular}{l l r r r}
\multicolumn{5}{c}{Triangulum VII.}\\[5pt]
Hauselberg&. . . . . .&\(34^0\)&\(25'\)&\(46''752\)\\
Wulfsode&. . . . . . .&109&38&36,566\\
Wilsede&. . . . . . . .&35&55&37,227
\end{tabular}
\end{center}

Aderunt itaque septem aequationes conditionales secundi generis (aequationes primi generis manifesto cessant), quas ut eruamus, computandi sunt ante omnia excessus sphaeroidici septem triangulorum. Ad hunc finem requiritur cognitio magnitudinis absolutae saltem unius lateris: latus inter puncta Wilsede et Wulfsode est 22377,94 metrorum. Hinc prodeunt excessus sphaeroidici triangulorum I\(.\,.\,.\;0''202\); II\(.\,.\,.\;2''442\); III\(.\,.\,.\;1''257\); IV\(.\,.\,.\;1''919\); V\(.\,.\,.\;1''957\); VI\(.\,.\,.\;0''321\); VII\(.\,.\,.\;1''295\).

Iam si directiones eo ordine, quo supra allatae indicibusque distinctae sunt, per \(v^{(0)},v^{(1)},v^{(2)},v^{(3)}\) etc. designantur, trianguli I anguli fiunt
\[
v^{(3)}-v^{(2)},\quad v^{(5)}-v^{(4)},\quad 360^0+v^{(7)}-v^{(10)}
\]
adeoque aequatio conditionalis prima
\[
-v^{(2)}+v^{(3)}-v^{(4)}+v^{(5)}+v^{(7)}-v^{(10)}+179^0\,59'\,59''798=0
\]
Perinde triangula reliqua sex alias suppeditant; sed levis attentio docebit, has septem aequationes non esse independentes, sed secundam identicam cum summa primae, quartae et sextae; nec non summam tertiae et quintae identicam cum summa quartae et septimae: quapropter secundam et quintam negligemus. Loco remanentium aequationum conditionalium in forma finita, adscribimus aequationes correspondentes e complexu (13), dum pro characteribus \(\varepsilon,\varepsilon'\) etc. his (0), (1), (2) etc. utimur:
\[
\begin{aligned}
-1''368&=-(2)+(3)-(4)+(5)+(7)-(10)\\
+1{,}773&=-(1)+(2)-(7)+(9)-(11)+(12)\\
+1{,}042&=-(0)+(2)-(7)+(8)+(14)-(17)\\
-0{,}813&=-(5)+(6)-(8)+(10)-(16)+(17)\\
-0{,}750&=-(8)+(9)-(11)+(13)-(15)+(17)
\end{aligned}
\]

Aequationes conditionales tertii generis \textit{octo} e triangulorum systemate peti possent, quum tum terna quatuor triangulorum I, II, IV, VI, tum terna ex his
